\section{Conclusion}
\label{sec:conclusion}

We have shown that self-conditioning in learned geometry models exhibits a fundamental asymmetry: rotation contracts toward a fixed point on SO(3) while translation diverges due to Sim(3) scale ambiguity. Exploiting this asymmetry via decoupled hybrid fusion---sourcing rotation from conditioned inference and translation from unconditioned inference---yields consistent improvements without the degradation of na\"ive iteration. Iterating twice achieves 27\% rotation improvement over the raw Pi3X baseline on a challenging wide-baseline, feature-adversarial calibration task, at only 3$\times$ the compute cost and with no external calibration data required.

Two findings merit emphasis. First, the contraction rate ($\lambda \approx 0.8$) and resulting accuracy floor ($\sim$3\degree{}) are properties of the model and scene, not the method---breaking this floor requires either stronger external priors or architectural changes to the conditioning mechanism. Second, Monte Carlo perturbation of the conditioning signal provides no measurable benefit: the posterior under self-conditioning is tightly unimodal ($<$1\degree{} angular standard deviation across samples), making deterministic iteration strictly preferable to stochastic ensembles. We believe the decoupled self-refinement framework generalises to any learned geometry model with a pose conditioning interface, and the contraction analysis provides a principled diagnostic for predicting when iterative refinement will help.
